\section{Related Work}

\subsection{Molecular property prediction benchmarks}

Public benchmarks are essential infrastructure for molecular machine learning. MoleculeNet brought together diverse molecular property prediction datasets and encouraged standardized comparisons across models, representations, and metrics \citep{wu2018moleculenet}. TDC further expanded this benchmark ecosystem by organizing therapeutics datasets into tasks and evaluation settings that better reflect drug discovery and development workflows \citep{huang2021tdc}. These resources make molecular machine-learning studies more comparable, but they also make an evaluation issue more visible: benchmark conclusions depend not only on datasets and models, but also on the split protocol used to define generalization.

Deep molecular learning studies further illustrate the importance of evaluation design. PotentialNet and related graph-based methods highlight the value of learned molecular representations for property prediction \citep{feinberg2018potentialnet}. Large-scale comparisons of learned representations and fingerprint-based methods have also shown that model performance can vary substantially across endpoints, datasets, and evaluation conditions \citep{yang2019learned}. The present study does not compete with these models. Instead, it examines whether the benchmark split changes the interpretation of model scores.

\subsection{Random and scaffold splitting}

Random splitting is simple and statistically convenient, but it can place structurally similar molecules in both training and test sets. In molecular property prediction, such similarity can inflate apparent generalization because the test set may contain near-neighbors of training compounds. Scaffold splitting is used as a stricter alternative, motivated by the idea that generalizing to unseen molecular frameworks is closer to the challenge faced in medicinal chemistry. The molecular framework concept underlying scaffold grouping is related to Bemis--Murcko scaffolds \citep{bemis1996properties}.

However, scaffold splitting does not eliminate all evaluation concerns. A scaffold split can still contain chemically similar molecules across different scaffolds, and recent work has argued that scaffold splits may overestimate virtual screening performance when scaffold dissimilarity does not fully capture chemical dissimilarity \citep{guo2024scaffold}. In addition, scaffold splitting may introduce target distribution shift if held-out scaffold groups have systematically different labels or property values. This means that a scaffold split can be chemically non-overlapping and statistically confounded at the same time.

\subsection{QSAR validation, applicability domain, and diagnostic evaluation}

The need for split diagnostics is consistent with long-standing QSAR and QSPR validation principles. External validation and careful interpretation of predictive performance have been emphasized as necessary conditions for reliable structure--property modelling \citep{tropsha2003earnest,gramatica2007principles}. Applicability-domain studies similarly stress that predictions are most reliable inside the chemical and descriptor space represented by the training data \citep{netzeva2005applicability,sahigara2012applicability}. Test-to-train chemical similarity is therefore not a secondary detail; it directly affects whether a benchmark resembles interpolation within a known chemical domain or extrapolation to a more distant region.

Many AIDD datasets are also imbalanced or distributionally uneven. Imbalanced benchmark studies such as ImDrug emphasize that class imbalance can distort model interpretation and metric choice in AI-aided drug discovery \citep{li2022imdrug}. Similar concerns apply to split-induced target shift: if train and test sets differ in label prevalence or property ranges, a performance difference may not be attributable to chemical generalization alone. This motivates diagnostic evaluation beyond a single performance number.

The present study follows this diagnostic view. Rather than treating random and scaffold splits as fixed benchmark choices, the analysis measures how split protocols change target distributions, scaffold composition, chemical similarity, model performance gaps, and model rankings. This positioning is intentionally different from a model leaderboard paper. The goal is not to show that one model is best, but to show that split-induced artifacts can change how benchmark results should be interpreted.
